% AST 623 — Asteroseismology in-class worksheet
% Goes with the interactive exercises on the course website:
%   Part 1: Examining the Power Spectrum
%   Part 2: Understanding Detectability
% Build:  latexmk -pdf worksheet.tex

\documentclass[11pt]{article}
\usepackage[margin=1in]{geometry}
\usepackage[T1]{fontenc}
\usepackage{lmodern}
\usepackage{amsmath}
\usepackage{siunitx}
\usepackage{enumitem}
\usepackage{xcolor}
\usepackage{needspace}
\usepackage[colorlinks=true, urlcolor=blue!60!black, linkcolor=black]{hyperref}

% ---- Notation -----------------------------------------------------------
\newcommand{\numax}{\ensuremath{\nu_\mathrm{max}}}
\newcommand{\dnu}{\ensuremath{\Delta\nu}}
\newcommand{\dnuzt}{\ensuremath{\delta\nu_{02}}}
\newcommand{\teff}{\ensuremath{T_\mathrm{eff}}}
\newcommand{\feh}{\ensuremath{[\mathrm{Fe/H}]}}
\newcommand{\uHz}{\si{\micro\hertz}}

% ---- Layout helpers -----------------------------------------------------
% Blank space for a written answer.
\newcommand{\answer}[1][3cm]{\par\vspace{#1}}
% A labelled blank to fill in, e.g. \fillin{\numax}{\uHz}
\newcommand{\fillin}[2]{#1 $=$ \underline{\hspace{3cm}} $\pm$ \underline{\hspace{2cm}}\ #2}
% Short instructions for what to do on the website.
\newcommand{\onsite}[1]{\par\noindent\textcolor{black!60}{\small\textit{On the website:} #1}\par\smallskip}

\setlist[enumerate,1]{label=\textbf{\arabic*.}, leftmargin=*, itemsep=0.5em}
\setlength{\parindent}{0pt}
\setlength{\parskip}{0.4em}

\title{\textbf{AST 623: Asteroseismology of Sun-like Stars}\\[0.3em]
       \large In-class worksheet}
\author{}
\date{}

\begin{document}
\maketitle
\vspace{-3em}

\noindent Name: \underline{\hspace{6cm}} \hfill Target star: KIC \underline{\hspace{3cm}}

\medskip
Work through this worksheet alongside the interactive exercises on the course website
(\url{https://TODO-github-pages-url/}). Pick your target star at the top of each exercise
page and use the same star throughout. Write down units and uncertainties for everything
you measure.

% =========================================================================
\section*{Part 1: Examining the Power Spectrum}

\needspace{12\baselineskip}\subsection*{The light curve}
\onsite{Panel 1. Hover to read values; drag a box to zoom; \emph{Reset view} zooms back out.}

\begin{enumerate}[series=ps]
  \item Zoom in until you can see individual one-minute points. What is the typical
        point-to-point scatter, in ppm?
        \answer[2cm]
  \item Solar-like oscillations have periods of a few minutes and amplitudes of a few ppm.
        Can you see them by eye? Why or why not?
        \answer
\end{enumerate}

\needspace{12\baselineskip}\subsection*{The power spectrum}
\onsite{Panel 2. Use the axes toggle (log--log / linear) and the smoothing box. Clicking pins
a frequency; the table under the plot reports the differences between pins
(\emph{Large frequency separation} mode) or within pairs of pins
(\emph{Small frequency separation} mode).}

\begin{enumerate}[resume=ps]
  \item In log--log view, find the broad, slowly declining background and the hump of
        oscillation power on top of it. Roughly what frequency range does the hump span?
        \answer[2cm]
  \item \textbf{\numax:} Smooth the spectrum with a width comparable to the mode spacing,
        around \SI{100}{\micro\hertz} or so. Then hover over the peak of the smoothed curve
        to estimate \numax.

        \fillin{\numax}{\uHz}
        \answer[1cm]
  \item Switch to linear axes and zoom into the hump. Click on a peak to pin its frequency.
        Pin a sequence of the tallest peaks. What spacing repeats?
        \answer[2.5cm]
  \item \textbf{\dnuzt:} Next to many strong peaks is a smaller peak just below it in
        frequency. Switch the pin table to \emph{Small frequency separation}, then pin a
        strong peak and its lower neighbour to measure the small separation. Do this for
        several orders and average.

        \fillin{\dnuzt}{\uHz}
        \answer[1cm]
\end{enumerate}

\needspace{12\baselineskip}\subsection*{The autocorrelation of the power spectrum}
\onsite{Panel 3. Set the frequency range and press \emph{Compute}. Click on the plot to pin
lags; the table lists the differences.}

\begin{enumerate}[resume=ps]
  \item \textbf{\dnu:} Compute the autocorrelation over a frequency range centred on the
        oscillation hump. Which lag gives the strongest peak? There is also a peak at about
        half that lag. What causes it?

        \fillin{\dnu}{\uHz}
        \answer[2.5cm]
  \item Now compute it over a range far from the hump (for example
        \SIrange{5000}{8000}{\micro\hertz}). What changes, and why?
        \answer
  \item Compare your \dnu\ from the autocorrelation with the spacing you measured by
        pinning peaks.
        \answer[2cm]
\end{enumerate}

\needspace{12\baselineskip}\subsection*{The \'echelle diagram}
\onsite{Panel 4. Set \dnu\ with the slider or the box. \emph{Show two \dnu\ widths} repeats the
diagram so ridges at the edges are not split. The $\varepsilon$ box shifts the horizontal axis
to $(\nu - \varepsilon\dnu) \bmod \dnu$.}

\begin{enumerate}[resume=ps]
  \item Enter your \dnu\ from the autocorrelation. Then nudge it up and down with the
        slider. How do the ridges tilt when \dnu\ is too small? Too large? Find the value
        that makes the strongest ridge most nearly vertical.
        \answer
  \item Identify the ridges of $\ell = 0$, 1 and 2 modes. (Hint: $\ell = 0$ and $\ell = 2$
        run close together; $\ell = 1$ sits roughly halfway between.) If a ridge is split
        across the left and right edges, tick \emph{Show two \dnu\ widths}.
        \answer
  \item \textbf{\dnuzt:} Hover over the $\ell = 0$ and $\ell = 2$ ridges in the same row to
        measure their horizontal separation. Does it change from the bottom of the diagram
        to the top? How does it compare with your answer to question~6?
        \answer
  \item Are the ridges perfectly straight? What might cause them to curve?
        \answer
\end{enumerate}

\needspace{12\baselineskip}\subsection*{Mass, radius and age from scaling relations}
\onsite{Panel 5. The star's \teff\ and \feh\ are filled in for you. Enter your seismic
measurements with their uncertainties.}

\begin{enumerate}[resume=ps]
  \item Enter your \numax, \dnu\ and \dnuzt\ with uncertainties. What mass, radius and age
        do you get?

        \begin{tabular}{@{}l@{\quad}l@{}}
          $M$     & $=$ \underline{\hspace{3cm}} $\pm$ \underline{\hspace{2cm}} $M_\odot$ \\[0.6em]
          $R$     & $=$ \underline{\hspace{3cm}} $\pm$ \underline{\hspace{2cm}} $R_\odot$ \\[0.6em]
          $\tau$  & $=$ \underline{\hspace{3cm}} $\pm$ \underline{\hspace{2cm}} Gyr
        \end{tabular}
        \answer[0.5cm]
  \item The age exponent on \dnu\ is about 9. If your \dnu\ were off by 1\%, how far off
        would the age be? Try it.
        \answer[2.5cm]
  \item Which single input dominates the uncertainty on the age? What would you need to
        measure better?
        \answer
\end{enumerate}

% =========================================================================
\newpage
\section*{Part 2: Understanding Detectability}
\onsite{The sliders at the top change the time span, the cadence and the added white noise.
The power spectrum, light curve and autocorrelation are recomputed each time. Pins in the
power spectrum and the autocorrelation stay put when you move the sliders.}

\needspace{12\baselineskip}\subsection*{Time span and frequency resolution}
\begin{enumerate}[series=det]
  \item The frequency resolution is about $1/T$, where $T$ is the time span. Shorten the
        data to 10 days. What is the resolution now? Can you still separate $\ell = 0$ and
        $\ell = 2$ modes?
        \answer
  \item Find the shortest time span for which the autocorrelation still gives a clear \dnu\
        peak. Which of \numax, \dnu\ and \dnuzt\ survives to the shortest span?
        \answer
\end{enumerate}

\needspace{12\baselineskip}\subsection*{Cadence and the Nyquist frequency}
\begin{enumerate}[resume=det]
  \item Increase the cadence to 2 minutes, then 5 minutes. Note the Nyquist frequency each
        time. Where does the oscillation power appear at 5-minute cadence, and why?
        \answer
  \item Averaging over a longer exposure also weakens signals at high frequency. Compare
        the height of the oscillation hump at 1, 2 and 5 minutes.
        \answer
  \item Kepler's long cadence was about 30 minutes. Could it detect oscillations in this
        star? What kind of star could it detect them in?
        \answer
\end{enumerate}

\needspace{12\baselineskip}\subsection*{Noise and signal-to-noise}
\begin{enumerate}[resume=det]
  \item Add noise until the oscillation hump disappears from the smoothed power spectrum.
        How much noise per point does that take?
        \answer[2cm]
  \item At that noise level, does the autocorrelation still show \dnu? Which method is more
        sensitive?
        \answer
  \item Now combine settings: with 30 days of data, how much noise can you tolerate before
        \dnu\ is lost? How does that compare with the full time span? How would you expect
        the detectability to scale with $T$?
        \answer[4cm]
\end{enumerate}

\end{document}
